% Curriculum Vitae Rodrigo Taborda -- two-page SHORT version (moderncv classic)
% Encoding: UTF-8. Full version: cv-rtaborda-english.tex
\documentclass[10pt,a4paper,sans]{moderncv}
\moderncvcolor{burgundy}
\moderncvstyle{classic}
\usepackage[utf8]{inputenc}
\usepackage[T1]{fontenc}
\usepackage[english]{babel}
\usepackage[scale=0.86]{geometry}
\setlength{\hintscolumnwidth}{2.6cm}

% ---------------------------------------------------------------------
% Datos personales
% ---------------------------------------------------------------------
    \name{Rodrigo}{Taborda}
    \title{Associate Professor · School of Management}
    \address{Calle 21 \# 1-20, School of Management}{Universidad de los Andes, Bogotá, Colombia}
    \email{r.taborda@uniandes.edu.co}
    \homepage{www.rodrigotaborda.com}
    \extrainfo{Nationality: Colombian}

% Publication entry: {year}{title}{coauthors}{source}
\newcommand{\pub}[4]{\cvitem{#1}{\textbf{#2}\ifx&#3&\else\ (with #3)\fi. #4}}

\begin{document}
\pagestyle{fancy}\fancyhf{}\renewcommand{\headrulewidth}{0pt}
\fancyfoot[C]{\addressfont\color{color2}Rodrigo Taborda --- Curriculum Vitae (short) --- \thepage}
\makecvtitle

    \section{Positions}
        \cventry{2013--}{Associate Professor}{School of Management, Universidad de los Andes}{}{}{}
        \cventry{2010--2013}{Lecturer}{School of Economics, Universidad del Rosario}{}{}{}
        \cventry{2010--2023}{Editor}{\emph{Revista de Economía del Rosario}}{}{}{}
        \cventry{2008--2010}{Consultant}{ACIL-Tasman, Canberra, Australia}{}{}{}
        \cventry{2002--2004}{Lecturer}{Universidad del Rosario}{}{}{}

    \section{Education}
        \cventry{2009}{Ph.D. Economics}{Australian National University}{}{}{Dissertation: \emph{Manufacturing industry performance after structural reforms. The case of Colombia}.}
        \cventry{2005}{M.A. Economics}{Australian National University}{}{}{}
        \cventry{2004}{M.A. Economics}{Universidad del Rosario}{}{}{}
        \cventry{2000}{B.A. Economics}{Universidad de Antioquia}{}{}{}

    \section{Publications}
        \subsection{Refereed journal articles}
            \pub{2027}{A test of auditor performance under the service paradigm: do emotional intelligence and national culture add to our understanding?}{Rosa González, Stephen B. Salter and Felipe Janica}{\emph{Journal of International Accounting, Auditing and Taxation}, forthcoming}
            \pub{2025}{Time poverty: an unintended consequence of women participation in farmers' associations}{Andrés Barrios and Ximena Rueda}{\emph{Social Indicators Research}, 178(1), 225--253}
            \pub{2023}{Fund flows, stock markets, and economic policy uncertainty: from the perspective of a CIVETS nation}{Joseph J. French and Philipp D. Schaberl}{\emph{Emerging Markets Finance and Trade}, 60(5), 967--982}
            \pub{2022}{Digitalization and e-government in the lives of urban migrants: evidence from Bogota}{Charles Martin, Sonia Camacho and Constantine Ruhe}{\emph{Policy \& Internet}, 14(2)}
            \pub{2022}{Savings and gender in Colombia: ANOVA, simple linear regression, and extreme observations}{Luz Orozco}{\emph{SAGE Research Methods Cases}}
            \pub{2019}{How does the training of subsistence entrepreneurs affect their hope levels and perceived goal-attainment?}{Andrés Barrios and Ezequiel Reficco}{\emph{Education + Training}, 61(7/8)}
            \pub{2019}{Family firms and financial performance. A meta-regression analysis}{Maximiliano González and Juan Idrobo}{\emph{Academia Revista Latinoamericana de Administración}, 32(3)}
            \pub{2018}{Disentangling the relationship between liquidity and returns in Latin America}{Joseph French}{\emph{Global Finance Journal}, 36}
            \pub{2017}{Stock liquidity and second blockholder as drivers of corporate value: evidence from Latin America}{Carlos Pombo}{\emph{International Review of Economics and Finance}, 51}
            \pub{2015}{Innovation and productivity in the Colombian service and manufacturing industries}{Juan Gallego and Luis Gutiérrez}{\emph{Emerging Markets Finance and Trade}, 51}
            \pub{2015}{Procedural transparency in Latin American central banks under inflation targeting schemes. A text analysis of the minutes of the boards of directors}{}{\emph{Ensayos sobre Política Económica}, Banco de la República, 33}
            \pub{2014}{Insurance use and economic growth in Latin America. Panel data evidence}{Ángela Concha}{\emph{Lecturas de Economía}, 81, July--December}
            \pub{2013}{Bias in economic news: the reporting of the nominal exchange rate behavior in Colombia}{}{\emph{Economía -- Journal of LACEA}, 14(1)}
            \pub{2011}{Manufacturing employment and wage differentials after structural adjustment reforms in Colombia}{Juan Guataquí}{\emph{Lecturas de Economía}, 48, January--June}
            \pub{2008}{Ownership and control in Colombian corporations}{Carlos Pombo and Luis Gutiérrez}{\emph{Quarterly Review of Economics and Finance}, 48(1)}
            \pub{2006}{Technical instruments and conceptual definitions: theoretical and empirical implications of the change on the definition of unemployment in Colombia}{Juan Guataquí}{\emph{Revista de Economía del Rosario}, 9(1)}
            \pub{2006}{Performance, efficiency and productivity in Colombia's power distribution system: effects of the 1994 regulatory reform}{Carlos Pombo}{\emph{Energy Economics}, 28(3)}
        \subsection{Books}
            \cvitem{}{\textbf{Ejercicios de econometría} (2016) and \textbf{Curso de macroeconomía} (2005), Editorial Universidad del Rosario.}

    \section{Teaching}
        \cvlistitem{29 different courses taught between 2002 and 2026 at four institutions: 19 at Universidad de los Andes and 8 at Universidad del Rosario (Colombia), and 2 outside Colombia (IESA, Venezuela; Universidad Americana, Paraguay), plus tutoring at the Australian National University.}
        \cvlistitem{14 undergraduate, 14 master and MBA, and 1 Ph.D., plus 2 executive education programs.}
        \cvlistitem{10 courses of original design and 4 online (Master in Analytics and Financial Management and online MBA).}
        \cvlistitem{At least 133 course-terms in total, 16 of them intensive intersemester courses.}
        \cvlistitem{Teaching in Spanish and English (International Master in Finance with Nova SBE) and co-teaching.}

    \section{Thesis supervision}
        \cvitem{}{Advisor of 2 Ph.D. dissertations in progress (Ph.D. in Management), 3 master theses (NOVA--Uniandes International Master in Finance; M.A. in Economics) and 2 undergraduate theses. Committee member or external reader of 6 undergraduate and 23 master theses.}

    \section{Translations}
        \cvitem{}{18 books translated from English into Spanish (World Bank, IDB, Universidad del Rosario, Alfaomega, Mayol), including \emph{Opportunism and goodwill} (Universidad del Rosario, forthcoming), \emph{Governing the excluded}, \emph{Gender and citizenship in transitional justice} (2024) and \emph{Autocracy and redistribution} (2021); and 3 technical reviews of translations (World Bank, IDB).}

    \section{Referee}
        \cvitem{Journals}{Academia Revista Latinoamericana de Administración, Economic Change and Restructuring, Economic Papers, Economic Systems, Emerging Markets Finance and Trade, Empirical Economics, Energy Economics, Energy Journal, Ensayos sobre Política Económica, Estudios Gerenciales, International Review of Economics and Finance, Journal of Business Research, Lecturas de Economía, Managerial Finance, Revista de Economía del Rosario, Sociedad y Economía, Stata Journal.}
        \cvitem{Projects}{Swiss National Science Foundation (SNSF), COLCIENCIAS, Banco de la República (FPIT) and Colombian universities. Jury member, Otto de Greiff Prize (2015, 2025).}

    \section{Awards, languages and software}
        \cvitem{2004}{Banco de la República (Colombia's central bank) scholarship for Ph.D. studies abroad.}
        \cvitem{2000}{Young researcher, COLCIENCIAS (Colombia's National Science Foundation).}
        \cvitem{Languages}{Spanish (native), English (C1).}
        \cvitem{Software}{Stata, R, EViews, Power BI, \LaTeX.}

\end{document}
